\section{野球ボールの空気力学的背景}
\label{sec:aerodynamics}

本章では，投球中の野球ボールに働く空気力学的な力を特徴づける無次元パラメータと，本プロジェクトが対象とする物理現象——縫い目に起因する非マグヌス力——について整理する。

\subsection{基本諸元と支配的な無次元パラメータ}
\label{sec:aero-params}

公式野球ボールの直径は円周の規定（9.00--9.25\,in）から72.9--74.9\,mmであり，本プロジェクトでは中央値 \(D = 0.0748\,\text{m}\) を既定値としてパラメータ化する。質量は既定 \(m = 0.145\,\text{kg}\)。表面には108個の二重縫い目からなる1本の連続曲線が球面上を二重らせん状に周回しており，その凸高さは実測で平均約0.79\,mm（\(D/92\) 程度）と報告されている\cite{kensrud2010thesis}。

\paragraph{2シーム・4シームについての注記}
野球ボールの縫い目は幾何学的には1種類（108個の二重縫い目からなる単一の閉曲線）しか存在しない。投げ方（グリップ）によって変わるのは，\emph{この固定された縫い目パターンに対する回転軸の向きと位相}である。回転軸をある向きに取ると1回転あたり4本の縫い目が気流を横切るように見え，別の向きでは2本に見える——これが「4シーム」「2シーム」という分類の正体であり，別々の幾何形状ではない。したがって本プロジェクトでは縫い目曲線を1つだけ実装し，投球条件として「回転軸の向き（チルト角・ヨー角）」と「縫い目パターンに対する初期位相角」を与える設計とする（\secref{seam-geometry}）。

代表投球速度は \(U = 20\)--45\,m/s，代表回転数は1500--3500\,rpmである。これらから定まる無次元パラメータのうち，本問題の物理を支配する3つを挙げる。

\begin{itemize}
	\item \textbf{レイノルズ数} \(Re = UD/\nu\)。代表条件（\(U=40\,\text{m/s}\)，\(\nu_{air}\approx1.5\times10^{-5}\,\text{m}^2/\text{s}\)）で \(Re\approx2.0\times10^5\)。文献では投球速度域で \(Re\approx1.1\text{--}2.4\times10^5\) と報告される\cite{nathan2008,kensrud2010thesis}。この値は滑面球の臨界レイノルズ数（ドラッグクライシス）が生じる遷移域に重なるため，境界層の層流--乱流遷移の位置が抗力に敏感に効く。
	\item \textbf{スピンパラメータ} \(S = \omega r/U\)（\(\omega\)：角速度，\(r\)：半径）。実投球では \(S\approx0.09\text{--}0.6\) 程度。
	\item \textbf{マッハ数} \(Ma = U/c\)。\(0.15\) 以下では圧縮性の影響は小さいが，格子ボルツマン法は元来弱圧縮性モデルであるため，格子単位系の設計において格子マッハ数を明示的に管理する必要がある（\secref{lattice-units}）。
\end{itemize}

\subsection{抗力とドラッグクライシス}

滑らかな球では，境界層が層流のまま剥離する亜臨界域で \(C_D\approx0.5\) 程度と大きいが，\(Re\) が臨界値を超えると境界層が乱流に遷移して剥離点が後方へ移動し，\(C_D\) が急減する（drag crisis）。

野球ボールでは縫い目が人工的な乱流トリップとして働くため，滑面球のような急峻なドラッグクライシスは示さず，縫い目の向きに依存してなだらかかつ非対称に \(C_D\) が変化することが実験的に確認されている\cite{kensrud2010thesis,kensrud2017study}。したがって本シミュレーションでは，滑面球の近似ではなく\emph{縫い目形状を含めた表面ジオメトリで境界層遷移・剥離を直接解像する}方針をとる。

後述するように，本プロジェクトの代表条件では縫い目の凸高さと層流境界層厚がほぼ一致する（\secref{resolution-budget}）。粗さ要素が境界層と同程度の高さを持つことは，それが遷移を誘発するトリップとして機能する条件そのものであり，縫い目を解像することと境界層を解像することは実質的に1つの要求である。

\subsection{マグヌス力と非マグヌス力（縫い目由来の横力）}
\label{sec:magnus}

回転する球のまわりでは，回転方向側の境界層が加速されて剥離が遅れ，反対側では剥離が早まることで循環が生じ，進行方向に垂直な力（マグヌス力）が発生する：
\begin{align}
	F_L = \frac{1}{2}\rho U^2 A\, C_L(S, Re, \text{seam orientation}).
	\label{eq:magnus}
\end{align}
実験的には \(S<0.4\) 程度の範囲で \(C_L\approx S\)（傾き1の線形関係）が良い近似となる\cite{nathan2008,lyu2011dragspin}が，この関係は滑面球の理想解に対する近似である。実際の野球ボールでは，縫い目の向きによって主流に対して非対称な圧力分布が生じ，マグヌス面（回転軸に直交する平面）から外れた方向の力——\textbf{非マグヌス力（横力）}——が発生することが知られている。これがスライダーやジャイロボール，スウィーパーが「教科書どおりの向きに曲がらない」現象の主要因である。

\textbf{本コードの設計要件}は，簡略化されたマグヌス力の解析式を用いるのではなく，縫い目形状を含む球体表面を境界条件として直接解像し，抗力係数 \(C_D\)・マグヌス方向の揚力係数 \(C_L\)・非マグヌス方向の横力係数 \(C_{side}\) をすべて流体場から直接算出することである。力の分解方法は\secref{coupling-loop}で述べる。

\subsection{運動方程式}
\label{sec:eom}

質量中心の並進運動：
\begin{align}
	m\frac{d\vec V}{dt} = \vec F_{drag} + \vec F_{Magnus} + \vec F_{side} + m\vec g.
	\label{eq:translation}
\end{align}
回転運動（角速度減衰，ジャイロ効果を含む一般形）：
\begin{align}
	I\frac{d\vec\omega}{dt} = \vec T_{aero}.
	\label{eq:rotation}
\end{align}
投球中の回転数減衰は数\%程度と小さいが，トルクは流体場から評価できる設計とする。本プロジェクトの密結合方式（\secref{coupling}）では，この運動方程式は係数テーブルではなく，CFDから毎ステップ得られる瞬時の力・トルクを直接使って積分する。
